\chapter{Background}\label{ch:background}

\section{Bike-Share Schemes}
Public bike-share schemes have spread to more than a thousand cities since the mid-2000s. Reviews of the research agree on two points \cite{fishman2013bike,shaheen2010bikesharing}. First, use depends heavily on how close stations are to where people start and end their journeys. Second, most trips replace walking or public transport rather than driving, so the effect on car travel is real but smaller than scheme operators tend to claim.

\section{What Influences Use}
A study of the scheme in Washington, DC, found that trips rose with temperature up to about 30 degrees and fell sharply in rain, and that stations near public transport were used more \cite{gebhart2014impact}. University schemes add a strong calendar effect, because most members leave during vacations.

\section{Lessons for This Pilot}
Three lessons shaped the design of the pilot:
\begin{itemize}
  \item stations were placed at the two ends of the most common walking journeys, not spread evenly;
  \item one station was placed at the railway station, to connect with public transport;
  \item the survey asked directly which mode each trip replaced, so that the effect on car travel would not have to be inferred.
\end{itemize}
